\chapter{D/E 题部分分策略与高级算法入口}\label{ch:08}

\section{D/E 题定位}\label{sec:08-1}

D/E 题可以先找小数据、特殊结构，或自己练过的做法。能独立完成一个子任务，就先把它写好；暂时不会的优化留到后面。

\begin{itemize}
  \item D 题先看子任务，按会写的做法估算能覆盖哪些数据。
  \item E 题也可以找小数据或特殊性质，但要给前面会做的题留够时间。
  \item 先完成选定子任务，再考虑扩大适用范围。
  \item 思考很久仍没进展时，可以先保存思路，转做其他题。
\end{itemize}

如果已经花了十几二十分钟仍没有完整思路，可以重新看子任务表，或者暂时换题。

\section{D/E 题读题流程}\label{sec:08-2}

\begin{enumerate}
  \item 先看有没有子任务或部分分说明。
  \item 抄下所有数据范围，尤其是小数据档。
  \item 找特殊性质：链、树、DAG、无修改、无查询、边权相同、参数很小。
  \item 写出最暴力做法。
  \item 估算暴力能拿哪些数据点。
  \item 如果会优化，再逐步优化。
\end{enumerate}

\section{部分分决策树}\label{sec:08-3}

\begin{longtable}{p{4.5cm}p{5cm}p{4.5cm}}
\toprule
题面或数据特征 & 可写做法 & 目标 \\
\midrule
\endfirsthead
\toprule
题面或数据特征 & 可写做法 & 目标 \\
\midrule
\endhead
$n \le 20$ & DFS、枚举子集、状态压缩 & 小数据分 \\
$n \le 100$ & $O(n^3)$、Floyd、区间 DP & 小中数据 \\
$n \le 1000$ & $O(n^2)$ & 中小数据 \\
图是树 & DFS、树形 DP、LCA 入门 & 特殊性质分 \\
图是链 & 转成数组问题 & 特殊性质分 \\
边权全为 1 & BFS & 替代 Dijkstra \\
没有修改 & 前缀和、静态预处理 & 静态数据分 \\
查询很少 & 每次暴力 & 小 $q$ 分 \\
修改很少 & 修改后重算 & 小修改分 \\
参数 $k$ 很小 & 状压、枚举 $k$ & 参数小分 \\
求阈值且可行性单调 & 二分答案（check 须正确） & 进一步优化 \\
\bottomrule
\end{longtable}

\section{暴力模板}\label{sec:08-4}

\subsection{枚举所有点对}

\begin{lstlisting}
long long ans = 0;
for (int i = 0; i < n; i++) {
    for (int j = i + 1; j < n; j++) {
        // check pair (i, j)
    }
}
\end{lstlisting}

适用：$n \le 5000$ 需谨慎，$n \le 1000$ 通常可试。

\subsection{枚举所有区间}

\begin{lstlisting}
for (int l = 0; l < n; l++) {
    long long sum = 0;
    for (int r = l; r < n; r++) {
        sum += a[r];
        // process interval [l, r]
    }
}
\end{lstlisting}

适用：$n \le 5000$ 以内看时限；如果 $n=10^5$，只能拿小数据分。

\subsection{DFS 枚举选择}

\begin{lstlisting}
int n;
vector<int> choose;

void dfs(int idx) {
    if (idx == n) {
        // evaluate choose
        return;
    }

    // not choose idx
    dfs(idx + 1);

    // choose idx
    choose.push_back(idx);
    dfs(idx + 1);
    choose.pop_back();
}
\end{lstlisting}

适用：$n \le 20$ 左右。

\section{小数据常用算法}\label{sec:08-5}

\subsection{Floyd}

点数小且问任意两点最短路：

\begin{lstlisting}
for (int k = 1; k <= n; k++) {
    for (int i = 1; i <= n; i++) {
        for (int j = 1; j <= n; j++) {
            if (dist[i][k] == INF || dist[k][j] == INF) continue;
            dist[i][j] = min(dist[i][j], dist[i][k] + dist[k][j]);
        }
    }
}
\end{lstlisting}

\subsection{全排列}

\begin{lstlisting}
vector<int> p(n);
iota(p.begin(), p.end(), 0);
do {
    // evaluate permutation
} while (next_permutation(p.begin(), p.end()));
\end{lstlisting}

适用：$n \le 9$ 或 $10$ 左右。

\subsection{状压枚举子集}

\begin{lstlisting}
for (int mask = 0; mask < (1 << n); mask++) {
    for (int i = 0; i < n; i++) {
        if (mask & (1 << i)) {
            // i is selected
        }
    }
}
\end{lstlisting}

适用：$n \le 20$ 左右。若 $n \ge 31$，注意 \Code{1LL << n} 和复杂度。

\section{特殊性质策略}\label{sec:08-6}

\subsection{图是链}

如果图是一条链，很多图问题可以转成数组问题：

\begin{itemize}
  \item 距离变成下标差或前缀和。
  \item 子树变成连续区间。
  \item 路径查询变成区间查询。
\end{itemize}

\subsection{图是树}

树比一般图简单：

\begin{itemize}
  \item 没有环。
  \item 两点之间路径唯一。
  \item 可用 DFS 统计子树。
  \item 可用树形 DP。
\end{itemize}

\subsection{边权相同}

如果所有边权都是 1：

\begin{itemize}
  \item 最短路用 BFS。
  \item 不需要 Dijkstra。
\end{itemize}

\subsection{无修改}

如果没有修改，只有查询：

\begin{itemize}
  \item 优先预处理。
  \item 区间问题考虑前缀和。
  \item 图问题考虑一次 BFS/Dijkstra 后回答多次查询。
\end{itemize}

\section{按题型拿部分分}\label{sec:08-7}

\subsection{图论题}

\begin{longtable}{p{4cm}p{5cm}p{5cm}}
\toprule
不会满分时 & 可写做法 & 可能覆盖 \\
\midrule
\endfirsthead
\toprule
不会满分时 & 可写做法 & 可能覆盖 \\
\midrule
\endhead
不会复杂最短路 & 小图 Floyd & $n$ 小的数据 \\
不会带权最短路 & 无权 BFS & 边权全 1 的数据 \\
不会动态图 & 每次修改后重跑 BFS/DFS & 操作数小的数据 \\
不会强连通 & 小图逐点搜索，判断双向可达 & $O(n(n+m))$，单次 DFS 不能代替强连通 \\
不会树上高级算法 & 从每个查询端点暴力爬/DFS & $n,q$ 小的数据 \\
\bottomrule
\end{longtable}

\subsection{区间数据结构题}

\begin{longtable}{p{4cm}p{5cm}p{5cm}}
\toprule
不会满分时 & 可写做法 & 可能覆盖 \\
\midrule
\endfirsthead
\toprule
不会满分时 & 可写做法 & 可能覆盖 \\
\midrule
\endhead
不会线段树 & 每次暴力扫描区间 & 小 $nq$ 数据 \\
只有查询无修改 & 前缀和 & 静态数据 \\
只有区间修改最后输出 & 差分 & 离线数据 \\
单点修改区间查询 & 树状数组入口 & 中等数据 \\
区间修改单点查询 & 差分思想 / 树状数组 & 中等数据 \\
\bottomrule
\end{longtable}

\subsection{DP 题}

\begin{longtable}{p{4cm}p{5cm}p{5cm}}
\toprule
不会满分时 & 可写做法 & 可能覆盖 \\
\midrule
\endfirsthead
\toprule
不会满分时 & 可写做法 & 可能覆盖 \\
\midrule
\endhead
状态想不清 & DFS 暴力枚举 & $n$ 小的数据 \\
背包优化不会 & 二维 DP & 小容量数据 \\
状态压缩不会 & 回溯搜索 & $n$ 小的数据 \\
转移太复杂 & 只处理部分状态或特殊参数 & 特殊性质数据 \\
\bottomrule
\end{longtable}

\subsection{字符串题}

\begin{longtable}{p{4cm}p{5cm}p{5cm}}
\toprule
不会满分时 & 可写做法 & 可能覆盖 \\
\midrule
\endfirsthead
\toprule
不会满分时 & 可写做法 & 可能覆盖 \\
\midrule
\endhead
不会 KMP & \Code{find} 或暴力匹配 & 短字符串数据 \\
不会 Trie & map 存完整字符串 & 字符串数量较少 \\
不会复杂解析 & split 简单格式 & 无嵌套数据 \\
不会哈希 & 直接比较字符串 & 小数据 \\
\bottomrule
\end{longtable}

\section{树状数组入口}\label{sec:08-8}

\subsection{什么题会用到}

\begin{itemize}
  \item 单点修改，前缀和或区间和查询。
  \item 需要动态维护数组和。
  \item 数据范围 $n,q \le 10^5$。
\end{itemize}

\subsection{树状数组模板}

\begin{lstlisting}
struct BIT {
    int n;
    vector<long long> tree;

    BIT(int n) : n(n), tree(n + 1, 0) {}

    void add(int idx, long long val) {
        assert(1 <= idx && idx <= n); // idx=0 would never advance
        for (; idx <= n; idx += idx & -idx) {
            tree[idx] += val;
        }
    }

    long long sumPrefix(int idx) {
        assert(0 <= idx && idx <= n);
        long long res = 0;
        for (; idx > 0; idx -= idx & -idx) {
            res += tree[idx];
        }
        return res;
    }

    long long rangeSum(int l, int r) {
        return sumPrefix(r) - sumPrefix(l - 1);
    }
};
\end{lstlisting}

\Warn{树状数组一般用 1 下标。}

\section{线段树入口}\label{sec:08-9}

\subsection{什么题会用到}

\begin{itemize}
  \item 区间查询和区间修改交替出现。
  \item 需要维护区间和、区间最大值、区间最小值。
  \item $n,q \le 10^5$，暴力扫描会超时。
\end{itemize}

\subsection{区间和线段树框架}

\begin{lstlisting}
struct SegTree {
    int n;
    vector<long long> tree, lazy;

    SegTree(int n) : n(n), tree(4 * n + 4), lazy(4 * n + 4) {}

    void apply(int node, int l, int r, long long val) {
        tree[node] += val * (r - l + 1);
        lazy[node] += val;
    }

    void push(int node, int l, int r) {
        if (lazy[node] == 0 || l == r) return;
        int mid = (l + r) / 2;
        apply(node * 2, l, mid, lazy[node]);
        apply(node * 2 + 1, mid + 1, r, lazy[node]);
        lazy[node] = 0;
    }

    void rangeAdd(int node, int l, int r, int ql, int qr, long long val) {
        if (ql <= l && r <= qr) {
            apply(node, l, r, val);
            return;
        }
        push(node, l, r);
        int mid = (l + r) / 2;
        if (ql <= mid) rangeAdd(node * 2, l, mid, ql, qr, val);
        if (qr > mid) rangeAdd(node * 2 + 1, mid + 1, r, ql, qr, val);
        tree[node] = tree[node * 2] + tree[node * 2 + 1];
    }

    long long query(int node, int l, int r, int ql, int qr) {
        if (ql <= l && r <= qr) return tree[node];
        push(node, l, r);
        int mid = (l + r) / 2;
        long long res = 0;
        if (ql <= mid) res += query(node * 2, l, mid, ql, qr);
        if (qr > mid) res += query(node * 2 + 1, mid + 1, r, ql, qr);
        return res;
    }
};
\end{lstlisting}

调用：

\begin{lstlisting}
// Require n >= 1; initial array is all zero.
SegTree seg(n);
// For each initial a[i], call rangeAdd(1,1,n,i,i,a[i]).
seg.rangeAdd(1, 1, n, l, r, v);
long long ans = seg.query(1, 1, n, l, r);
\end{lstlisting}

\subsection{不会线段树时怎么办}

\begin{itemize}
  \item 如果没有修改，用前缀和。
  \item 如果只有区间修改、最后查询，用差分。
  \item 如果是单点修改、区间查询，用树状数组。
  \item 如果数据小，直接暴力扫描。
\end{itemize}

\section{Kruskal 最小生成树入口}\label{sec:08-10}

\subsection{什么题会用到}

\begin{itemize}
  \item 连接所有点。
  \item 总代价最小。
  \item 无向图。
  \item 选择若干边让图连通。
\end{itemize}

\subsection{Kruskal 模板}

\begin{lstlisting}
struct Edge {
    int u, v;
    long long w;
};

sort(edges.begin(), edges.end(), [](const Edge& a, const Edge& b) {
    return a.w < b.w;
});

DSU dsu(n);
long long ans = 0;
int cnt = 0;
for (auto e : edges) {
    if (dsu.unite(e.u, e.v)) {
        ans += e.w;
        cnt++;
    }
}

if (cnt == n - 1) cout << ans << '\n';
else cout << "Impossible\n";
\end{lstlisting}

\section{树上问题入口}\label{sec:08-11}

\subsection{什么题会用到}

\begin{itemize}
  \item 给定一棵树。
  \item 查询两点路径。
  \item 查询祖先、最近公共祖先。
  \item 对树上路径加值或统计经过次数。
\end{itemize}

\subsection{树上暴力拿分}

如果不会 LCA，且 $n,q$ 小，可以每次 DFS 找路径。

\begin{lstlisting}
bool dfsPath(int u, int parent, int target, vector<int>& path) {
    if (u == target) {
        path.push_back(u);
        return true;
    }
    for (int v : g[u]) {
        if (v == parent) continue;
        if (dfsPath(v, u, target, path)) {
            path.push_back(u);
            return true;
        }
    }
    return false;
}
\end{lstlisting}

\subsection{LCA 倍增入口}

\begin{lstlisting}
int LOG; // compute from n before allocating up
vector<vector<int>> up;
vector<int> depth;
vector<vector<int>> g;

void dfsLca(int u, int p) {
    up[0][u] = p;
    for (int k = 1; k < LOG; k++) {
        up[k][u] = up[k - 1][up[k - 1][u]];
    }
    for (int v : g[u]) {
        if (v == p) continue;
        depth[v] = depth[u] + 1;
        dfsLca(v, u);
    }
}

int lca(int a, int b) {
    if (depth[a] < depth[b]) swap(a, b);
    int diff = depth[a] - depth[b];
    for (int k = 0; k < LOG; k++) {
        if (diff & (1 << k)) a = up[k][a];
    }
    if (a == b) return a;
    for (int k = LOG - 1; k >= 0; k--) {
        if (up[k][a] != up[k][b]) {
            a = up[k][a];
            b = up[k][b];
        }
    }
    return up[0][a];
}
\end{lstlisting}

先读入 n 并将无向树边存入 g。下面根为 1，根父亲为自身；深链应改用 \Tref{209} 的无递归遍历。初始化：

\begin{lstlisting}
LOG = 1;
while ((1LL << LOG) <= n) LOG++;
up.assign(LOG, vector<int>(n + 1));
depth.assign(n + 1, 0);
dfsLca(1, 1);
\end{lstlisting}

\subsection{树上差分入口}

多次给路径 $(u,v)$ 加 1，最后统计每个点或边被经过次数。

点差分采用根的父亲为自身的约定，经过根的路径不能再减一次根；以下函数与 LCA 共用 g、up。处理全部路径后调用 collect(1,1)。

\begin{lstlisting}
vector<long long> diff(n + 1, 0);

void addPath(int u, int v) {
    int w = lca(u, v);
    diff[u]++;
    diff[v]++;
    diff[w]--;
    if (up[0][w] != w) diff[up[0][w]]--; // root is its own parent
}

void collect(int u, int p) {
    for (int v : g[u]) {
        if (v == p) continue;
        collect(v, u);
        diff[u] += diff[v];
    }
}
\end{lstlisting}

点差分和边差分的公式不同。拿不准时，可先逐条路径更新，用小数据验证；是否能得部分分要看子任务范围。

\section{状态压缩 DP 入口}\label{sec:08-12}

\subsection{什么题会用到}

\begin{itemize}
  \item $n$ 很小，例如 $n \le 20$。
  \item 每个元素只有选或不选。
  \item 状态可以用二进制集合表示。
  \item 旅行、分组、覆盖、集合最优值。
\end{itemize}

\subsection{集合 DP 基本形状}

\begin{lstlisting}
int N = 1 << n;
vector<int> dp(N, INF);
dp[0] = 0;

for (int mask = 0; mask < N; mask++) {
    for (int i = 0; i < n; i++) {
        if ((mask & (1 << i)) == 0) {
            int nmask = mask | (1 << i);
            dp[nmask] = min(dp[nmask], dp[mask] + cost[i]);
        }
    }
}
\end{lstlisting}

\subsection{状压易错点}

\begin{itemize}
  \item $1 << n$ 当 $n \ge 31$ 会溢出，要用 \Code{1LL << n}。
  \item 状态数是 $2^n$，$n=25$ 已经很大。
  \item 数组大小和内存要提前估算。
\end{itemize}

\section{KMP 字符串匹配入口}\label{sec:08-13}

\subsection{什么题会用到}

\begin{itemize}
  \item 长文本中多次查找模式串。
  \item 字符串长度很大，暴力匹配会超时。
  \item 需要统计模式串出现次数。
\end{itemize}

\subsection{KMP 模板}

\begin{lstlisting}
vector<int> prefixFunction(const string& p) {
    int n = (int)p.size();
    vector<int> pi(n, 0);
    for (int i = 1; i < n; i++) {
        int j = pi[i - 1];
        while (j > 0 && p[i] != p[j]) j = pi[j - 1];
        if (p[i] == p[j]) j++;
        pi[i] = j;
    }
    return pi;
}

int countMatch(const string& text, const string& p) {
    if (p.empty()) return (int)text.size() + 1; // empty pattern: all gaps
    vector<int> pi = prefixFunction(p);
    int j = 0;
    int cnt = 0;
    for (char c : text) {
        while (j > 0 && c != p[j]) j = pi[j - 1];
        if (c == p[j]) j++;
        if (j == (int)p.size()) {
            cnt++;
            j = pi[j - 1];
        }
    }
    return cnt;
}
\end{lstlisting}

\subsection{不会 KMP 时}

\begin{itemize}
  \item 字符串短时用暴力匹配。
  \item C++ 可用 \Code{s.find(t)} 拿部分数据。
  \item 多模式串匹配可能要 Trie 或 AC 自动机，但新手优先拿小数据。
\end{itemize}

\section{Trie 字典树入口}\label{sec:08-14}

\subsection{什么题会用到}

\begin{itemize}
  \item 很多字符串插入和查询。
  \item 前缀查询。
  \item 字典、单词集合。
\end{itemize}

\subsection{小写字母 Trie}

\begin{lstlisting}
struct Trie {
    struct Node {
        int child[26];
        bool end = false;
        Node() {
            memset(child, -1, sizeof child);
        }
    };

    vector<Node> tr;

    Trie() {
        tr.push_back(Node());
    }

    void insert(const string& s) {
        int u = 0;
        for (char c : s) {
            int x = c - 'a';
            if (tr[u].child[x] == -1) {
                tr[u].child[x] = (int)tr.size();
                tr.push_back(Node());
            }
            u = tr[u].child[x];
        }
        tr[u].end = true;
    }

    bool find(const string& s) {
        int u = 0;
        for (char c : s) {
            int x = c - 'a';
            if (tr[u].child[x] == -1) return false;
            u = tr[u].child[x];
        }
        return tr[u].end;
    }
};
\end{lstlisting}

\subsection{不会 Trie 时}

\begin{itemize}
  \item 只查完整字符串，用 \Code{set<string>} 或 \Code{unordered\_set<string>}。
  \item 数据小，用逐个字符串比较。
  \item 前缀查询可以用排序后 \Code{lower\_bound} 尝试部分分。
\end{itemize}

\section{D/E 考场时间策略}\label{sec:08-15}

\begin{longtable}{p{3cm}p{8cm}p{3cm}}
\toprule
剩余时间 & 建议动作 & 目标 \\
\midrule
\endfirsthead
\toprule
剩余时间 & 建议动作 & 目标 \\
\midrule
\endhead
60 分钟以上 & 先思考满分算法，再写部分分保底 & 进一步优化 \\
30--60 分钟 & 先写暴力或特殊性质，再看能否优化 & 稳定拿分 \\
15--30 分钟 & 只写最确定的小数据做法 & 不空题 \\
15 分钟以内 & 完成一个明确子任务的正确实现，清理输出 & 完成可做的部分 \\
\bottomrule
\end{longtable}

\section{常见高级算法入口}\label{sec:08-16}

\begin{longtable}{p{4cm}p{5cm}p{5cm}}
\toprule
题面信号 & 可能算法 & 先翻哪里 \\
\midrule
\endfirsthead
\toprule
题面信号 & 可能算法 & 先翻哪里 \\
\midrule
\endhead
最大化最小值、最小化最大值 & 二分答案 & B 题二分答案 \\
区间修改查询交替 & 线段树 / 树状数组 & \Tref{401} / \Tref{402} \\
树上路径、祖先 & LCA / 树上差分 & \Tref{208}、本章树上差分 \\
强连通、缩点 & Tarjan & 本书未收录 Tarjan；小图可枚举可达性 \\
最小连接代价 & 最小生成树 & Kruskal / Prim \\
网络最短距离 & Dijkstra / Floyd & B 题图论入口 \\
容量和价值 & 背包 DP & B 题 DP 入口 \\
集合状态很小 & 状态压缩 DP & 状压入口 \\
字符串匹配很多次 & KMP / Trie & \Tref{502} / \Tref{503} \\
\bottomrule
\end{longtable}

\section{高级算法识别速查}\label{sec:08-17}

\begin{longtable}{p{4cm}p{5cm}p{5cm}}
\toprule
关键词 & 优先想到 & 不会时的替代 \\
\midrule
\endfirsthead
\toprule
关键词 & 优先想到 & 不会时的替代 \\
\midrule
\endhead
区间加、区间查询 & 线段树 & 差分、树状数组、暴力 \\
单点修改、区间和 & 树状数组 & 前缀和静态版、暴力 \\
树上路径查询 & LCA & 每次 DFS 找路径 \\
多次树上路径加 & 树上差分 & 暴力路径加 \\
连接所有点最小代价 & Kruskal & 枚举小图、Prim 小数据 \\
集合状态很小 & 状压 DP & DFS 枚举 \\
长串匹配短串 & KMP & \Code{find} / 暴力 \\
前缀查询很多 & Trie & set / 排序后二分 \\
最大化最小值 & 二分答案 & 枚举答案小范围 \\
任务依赖 & 拓扑排序 & 小图 DFS 检查 \\
\bottomrule
\end{longtable}

\section{D/E 模板练习建议}\label{sec:08-18}

\begin{longtable}{p{4cm}p{4cm}p{6cm}}
\toprule
模板 & 优先级 & 原因 \\
\midrule
\endfirsthead
\toprule
模板 & 优先级 & 原因 \\
\midrule
\endhead
二分答案 & 高 & B/D 都常见，代码短，提分明显。 \\
树状数组 & 高 & 区间题常见，比线段树简单。 \\
Dijkstra & 高 & 图论常见，B/D 都可能用。 \\
Kruskal & 中 & 最小生成树题专用，模板短。 \\
线段树 & 中高 & 强但代码长，建议练熟基础版。 \\
LCA & 中 & 树上题常见，但新手可先会暴力。 \\
状态压缩 DP & 中 & 看到 $n \le 20$ 很有用。 \\
KMP & 中低 & 字符串题有用，但 CSP 中不一定高频。 \\
Trie & 中低 & 前缀题有用，完整字符串可用 set 替代。 \\
树上差分 & 低到中 & 需要 LCA 基础，适合进一步优化。 \\
\bottomrule
\end{longtable}

\section{完整做法暂时不会时}\label{sec:08-19}

出现以下情况，应转为部分分：

\begin{itemize}
  \item 已经想了较久，仍没有可实现的做法。
  \item 会算法名字，但模板完全不会写。
  \item 写到一半发现状态定义不清。
  \item 样例都无法解释。
  \item A/B/C 还没稳。
\end{itemize}

\section{D/E 题提交前检查}\label{sec:08-20}

\begin{itemize}
  \item 是否至少能过样例。
  \item 是否能过自己造的小数据。
  \item 是否处理了最小输入。
  \item 如果是暴力，是否明确知道能拿哪些数据。
  \item 是否删除调试输出。
  \item 是否没有因为追求优化改坏暴力基础逻辑。
\end{itemize}
